\documentclass[11pt]{article}
\usepackage[margin=1in]{geometry}
\usepackage{graphicx}
\usepackage{amsmath, amssymb}
\usepackage{booktabs}
\usepackage{listings}
\usepackage{xcolor}
\usepackage{hyperref}
\usepackage{caption}
\usepackage{enumitem}

\definecolor{codebg}{RGB}{245,245,250}
\definecolor{cmdbg}{RGB}{40,42,54}
\definecolor{cmdfg}{RGB}{248,248,242}

\lstdefinestyle{py}{
    language=Python,
    basicstyle=\ttfamily\small,
    backgroundcolor=\color{codebg},
    keywordstyle=\bfseries\color{blue!70!black},
    stringstyle=\color{green!50!black},
    commentstyle=\itshape\color{gray},
    breaklines=true, frame=single, framesep=4pt, showstringspaces=false}

\lstdefinestyle{shell}{
    language=bash,
    basicstyle=\ttfamily\footnotesize\color{cmdfg},
    backgroundcolor=\color{cmdbg},
    breaklines=true, frame=single, framesep=4pt, showstringspaces=false}

\hypersetup{colorlinks=true, urlcolor=blue, linkcolor=black,
            citecolor=black, breaklinks=true}

\title{\Large\bfseries triloop \& picoacq\\[2pt]
       \large Analysis Code Reference\\[8pt]
       \normalsize Pipeline math, module-by-module walkthrough, and
       per-program description}
\author{}
\date{\today \\[6pt] \small Version 0.1.0}

\begin{document}
\maketitle

\tableofcontents
\bigskip\hrule\bigskip

%-----------------------------------------------------------------------
\section{Introduction}
%-----------------------------------------------------------------------

This document is a developer's reference for the analysis code in
\textsf{triloop} and the acquisition code in \textsf{picoacq}.  It
complements the \textit{triloop User Manual} (which describes how to
\emph{use} the tools) and the \textit{picoacq User Manual} (which
describes the hardware and the capture file format).  Here we focus on
\emph{what each Python module does}, in what order they are called by
the public entry points, and the mathematics each step implements.

\paragraph{Audience.} Reader is comfortable with Python, complex
baseband signal processing, and standard linear-algebra notation.

\paragraph{Conventions used in this document.} Bold lower-case
($\mathbf b$) denotes a vector; bold upper-case ($\mathsf{N}$) denotes
a matrix; $\hat{\mathbf k}$ is the unit propagation vector \emph{toward
the source}; angles are in degrees in code, radians in equations.

%-----------------------------------------------------------------------
\section{Pipeline overview}
%-----------------------------------------------------------------------

A capture flows through the toolchain in four stages:

\begin{enumerate}\itemsep=2pt
\item \textbf{Acquisition} (\textsf{picoacq}).  Hardware capture
      orchestrated by \texttt{picoacq.recorder.capture()}, which calls
      the PicoSDK wrapper \texttt{picoacq.ps5444a} or, in
      \texttt{--simulate} mode, \texttt{picoacq.simulator}.  Output:
      a triloop-format HDF5 file with float32 samples in volts plus a
      JSON \texttt{capture\_settings} dict carrying the chosen sample
      rate, voltage ranges, auto-range probe values, and (new in v0.1)
      the list of expected RF bands \texttt{rf\_bands\_hz}.
\item \textbf{Band extraction} (\texttt{triloop.bands},
      \texttt{triloop.extract}).  Read the HDF5 file with
      \texttt{triloop.io\_hdf5.read\_capture()}; identify the baseband
      alias position of each RF band via \texttt{picoacq.alias};
      mix-and-LPF each channel down to a complex baseband centred on
      that band; decimate to a workable IF rate.
\item \textbf{Polarimetry and direction finding}
      (\texttt{triloop.analyze}, \texttt{triloop.geometry},
      \texttt{triloop.stokes}, \texttt{triloop.direction},
      \texttt{triloop.beamform}).  Recover the lab-frame
      $\mathbf B$-vector from the three loops via the inverse of the
      geometry matrix $\mathsf{N}$, project onto the plane
      perpendicular to $\hat{\mathbf k}$, compute Stokes parameters
      and intensity, and (optionally) refine $\hat{\mathbf k}$ with
      a null-eigenvector / null-sweep search.
\item \textbf{Reporting} (\texttt{triloop.view},
      \texttt{triloop.multiband}, \texttt{triloop.browse}).  Static
      QC plots, multi-band comparison figures, and an interactive
      single-file HTML browser for visual inspection.
\end{enumerate}

The CLI tools (\texttt{triloop view}, \texttt{triloop analyze-multi},
\texttt{triloop browse}, \texttt{triloop analyze},
\texttt{triloop locate}, \texttt{triloop batch}, \texttt{triloop
summary}, \texttt{triloop simulate}, \texttt{triloop notebook}) wrap
these stages in different combinations.

%-----------------------------------------------------------------------
\section{Acquisition: \textsf{picoacq}}
%-----------------------------------------------------------------------

\subsection{\texttt{picoacq.cli}}

Click-based command-line interface with two subcommands:
\texttt{capture} (one-shot) and \texttt{monitor} (continuous).  Each
exposes the same configuration knobs as the underlying recorder; the
defaults reflect the multi-band undersampling workflow:

\begin{itemize}\itemsep=2pt
\item \texttt{--rate 15625000} (Hz)
\item \texttt{--duration 8.0} (s)
\item \texttt{--channels A,B,C,D}
\item \texttt{--range 2.0} (V; only used when \texttt{--no-auto-range})
\item \texttt{--auto-range/--no-auto-range} (default on)
\item \texttt{--headroom 3.0}
\item \texttt{--probe 0.1} (auto-range probe duration, s)
\item \texttt{--rf-bands "2.5e6,5e6,10e6,15e6,20e6"}
\end{itemize}

\subsection{\texttt{picoacq.ps5444a}}

A thin \texttt{ctypes} wrapper around the PicoSDK 5000a driver.  Five
functions:

\begin{itemize}\itemsep=2pt
\item \texttt{open\_5444(n\_channels=4)} -- context manager that opens
      the device, automatically picking 16-bit resolution for $\le 2$
      channels and 14-bit for 3--4 channels (5444D constraint).
\item \texttt{configure\_channel(ps, handle, ch, range\_volts,
      coupling)} -- enables a channel and sets its $\pm$range.
\item \texttt{get\_timebase(ps, handle, target\_rate)} -- searches the
      Pico timebase index space for the closest valid sample rate at
      or below the requested target.  Skips early indices that are
      invalid in 4-channel mode.
\item \texttt{acquire\_block(ps, handle, channels, n, tb, ranges)} --
      block-mode acquisition with explicit ctypes-wrapped arguments
      to avoid signature mismatches between picosdk versions.
\item \texttt{\_ensure\_pico\_dyld\_path()} -- macOS-only shim that
      injects the framework library path into
      \texttt{DYLD\_LIBRARY\_PATH} before importing \texttt{picosdk}.
\end{itemize}

\subsection{\texttt{picoacq.simulator}}

The no-hardware fallback.  Generates a 3-loop synthetic dataset with
configurable arrival direction, polarization, additive Gaussian noise,
and Faraday rotation rate.  Used by \texttt{--simulate} mode and as
the silent fallback when the SDK or hardware is unavailable.

\subsection{\texttt{picoacq.recorder}}

The orchestrator.  Top-level function:
\texttt{capture(output\_file, duration\_s, sample\_rate, channels,
ranges\_volts, coupling, loops\_config, simulate, sim\_kwargs,
auto\_range\_enabled, headroom, probe\_duration\_s, rf\_bands,
verbose)}.  Internally:

\begin{enumerate}\itemsep=2pt
\item \texttt{\_check\_buffer()} validates that
      $\mathrm{duration}\times\mathrm{rate}$ samples fit in the
      onboard memory (128\,MS/ch in 4-channel mode); raises
      \texttt{ValueError} with a clear maximum if not.
\item If auto-range is on (and \texttt{ranges\_volts} not supplied):
      open the scope at the maximum $\pm 20$\,V on every channel,
      run a brief probe capture (\texttt{auto\_range()}), measure
      peak/RMS in volts, and call \texttt{\_pick\_range()} to choose
      the smallest available Pico range whose full scale
      $\geq K\cdot$peak (default $K = 3$).
\item Run the science capture via \texttt{\_try\_real\_capture()};
      on any SDK/hardware failure, fall back transparently to
      \texttt{simulate\_capture()}.
\item Convert int16 ADC codes $\to$ float32 volts; write to HDF5 via
      \texttt{triloop.io\_hdf5.write\_capture()} with the chosen
      ranges, the auto-range probe results, and \texttt{rf\_bands\_hz}
      stored in \texttt{capture\_settings}.
\end{enumerate}

\subsection{\texttt{picoacq.alias}}
\label{sec:alias}

Pure-Python helper that maps RF frequencies to baseband alias positions
for a given sample rate.  An RF tone at $f_{\rm rf}$ folds to baseband as
\begin{equation}
f_{\rm bb} = \bigl|\,f_{\rm rf} - n\,f_s\,\bigr|,
\qquad n = \mathrm{round}(f_{\rm rf}/f_s),
\end{equation}
with Nyquist zone $z = \lfloor 2 f_{\rm rf}/f_s\rfloor + 1$.  Even
zones produce a frequency-inverted spectrum.  Three exports:

\begin{itemize}\itemsep=2pt
\item \texttt{Alias} dataclass with fields
      \texttt{(rf\_hz, baseband\_hz, nyquist\_zone, inverted)}.
\item \texttt{alias\_of(rf\_hz, sample\_rate\_hz)} -- single mapping.
\item \texttt{rf\_for\_baseband(f\_bb, rf\_bands, fs, tol)} --
      inverse lookup (given a detected baseband peak, return which
      RF band it most likely came from).
\end{itemize}

This module is the single source of truth for alias bookkeeping;
\texttt{triloop.bands} and \texttt{triloop.view}/\texttt{browse} all
import from it.

%-----------------------------------------------------------------------
\section{The capture file format}
\label{sec:fileformat-ref}
%-----------------------------------------------------------------------

Both packages read and write the same HDF5 v0.1 layout:

\begin{lstlisting}[style=shell]
/                           HDF5 root, @format_version = "0.1"
/metadata                   group, attributes:
    @sample_rate            float64, Hz
    @duration_s             float64
    @start_time_utc         ISO-8601
    @scope_model            string
    @scope_serial           string
    @capture_settings       JSON-encoded dict
    @loops_config           JSON-encoded dict
/channels
    /A   /B   /C   /D       float32 (or int16) datasets, length N
    each has:
        @gain_v_per_count   float64
        @offset_v           float64
/time                       optional float64 array, length N
\end{lstlisting}

\noindent\texttt{capture\_settings} fields recorded by
\texttt{picoacq.recorder.capture()}:
\begin{itemize}\itemsep=2pt
\item \texttt{requested\_duration\_s, requested\_sample\_rate}
\item \texttt{channels, ranges\_volts, coupling}
\item \texttt{used\_simulator, auto\_range\_enabled, headroom}
\item \texttt{auto\_range\_probe} (\texttt{peaks\_volts, rms\_volts,
      headroom, probe\_duration\_s}) -- only when probe ran
\item \texttt{picoscope\_ranges\_volts} -- the actual chosen ranges
\item \texttt{rf\_bands\_hz} -- list of expected RF frequencies in Hz
\item \texttt{truth} -- only present in simulator mode
\end{itemize}

\texttt{triloop.io\_hdf5} provides \texttt{read\_capture()} and
\texttt{write\_capture()}; both are short and tested, and any third
party can interoperate by writing the same layout.

%-----------------------------------------------------------------------
\section{Geometry and direction transforms}
%-----------------------------------------------------------------------

\subsection{\texttt{triloop.geometry}}

Operates on the loop-array geometry config dict
(\texttt{loops\_config}).  Five public functions:

\begin{itemize}\itemsep=2pt
\item \texttt{az\_el\_to\_khat(az\_deg, el\_deg)} returns
      $\hat{\mathbf k} = (\cos\theta\sin\varphi, \cos\theta\cos\varphi,
      \sin\theta)^\top$ in the local ENU frame.
\item \texttt{build\_N\_matrix(loops\_config)} returns the $3\times 3$
      matrix $\mathsf{N}$ whose $i$-th row is the $i$-th loop's normal
      $\hat{\mathbf n}_i$ (in ENU, scaled by per-loop gain).  This is
      the matrix that maps the lab-frame $\mathbf B$ to the
      loop-projected real signal $\mathbf b = \mathsf{N}\mathbf B$.
\item \texttt{perp\_projector(khat)} returns
      $\mathsf{P}_\perp = \mathsf{I} - \hat{\mathbf k}\hat{\mathbf k}^\top$,
      the projector onto the plane perpendicular to $\hat{\mathbf k}$.
\item \texttt{perp\_orthonormal\_basis(khat)} returns two unit vectors
      $(\hat{\mathbf p}, \hat{\mathbf q})$ that span that plane, with
      $\hat{\mathbf p}$ chosen along the local horizontal whenever
      possible, $\hat{\mathbf q} = \hat{\mathbf k}\times\hat{\mathbf p}$.
\item \texttt{LOOP\_NORMALS\_DEFAULT} and \texttt{\_normal\_from\_az\_el}
      support the default cube-vertex layout.
\end{itemize}

The default geometry is three loops on three faces of a cube whose
body diagonal is vertical, giving normal-tilts of $\arctan(1/\sqrt{2})
\approx 35.26^\circ$ above horizontal at azimuths $\{0, +120, -120\}^\circ$.

%-----------------------------------------------------------------------
\section{Band extraction and analysis}
%-----------------------------------------------------------------------

\subsection{\texttt{triloop.extract}}

Two functions covering the single-band path:

\begin{itemize}\itemsep=2pt
\item \texttt{extract\_complex\_baseband(t, b, f0, BW)} performs an
      FFT-based peak search around $f_0$ over a $\pm$\texttt{search\_BW}
      window, mixes the strongest peak to DC, brick-wall low-passes to
      $\pm$\texttt{BW}/2, and returns the complex baseband at the
      original sample rate, the detected peak frequency, and an
      FFT-domain SNR estimate.
\item \texttt{extract\_three\_loops(...)} calls the above on each loop,
      applies per-loop \texttt{phase\_offset\_deg} calibration, and
      returns a stacked $(3, N)$ complex array.
\end{itemize}

\subsection{\texttt{triloop.bands}}

Multi-band counterpart, used when the capture's
\texttt{rf\_bands\_hz} populates more than one entry.  Public
function: \texttt{extract\_bands(cap, rf\_bands, bw\_hz, decim\_rate\_hz, channels)}.
For every band:

\begin{enumerate}\itemsep=2pt
\item \texttt{picoacq.alias.alias\_of()} computes
      $f_{\rm bb}$, the Nyquist zone, and \texttt{inverted}.
\item Validate that the requested \texttt{bw\_hz} doesn't reach into
      the adjacent Nyquist edge or another band's alias.
\item Mix every selected channel down to $f_{\rm bb}$ via FFT
      (\texttt{\_extract\_one\_channel}), brick-wall filter to
      $\pm$\texttt{bw\_hz}/2, decimate to the smallest integer factor
      that yields $\geq$ \texttt{decim\_rate\_hz}, and conjugate the
      result if the band lives in an even (inverted) Nyquist zone.
\end{enumerate}

Returns a dict mapping \texttt{rf\_hz} $\to$ \texttt{BandExtraction}
dataclass with fields
\texttt{rf\_hz, baseband\_hz, nyquist\_zone, inverted, bw\_hz,
decim\_rate\_hz, t, z}, where \texttt{z} is a per-channel dict of
\texttt{complex64} arrays.

\subsection{\texttt{triloop.analyze}}

The polarimetry pipeline.  Two public entry points:

\begin{itemize}\itemsep=2pt
\item \texttt{analyze(t, B1, B2, B3, f0, BW, az\_deg, el\_deg, loops\_config)}
      -- starts from real loop signals, calls
      \texttt{extract\_three\_loops()}, then forwards to:
\item \texttt{analyze\_z\_loops(t, Z\_loops, f\_peak, snrs\_db, az\_deg, el\_deg, loops\_config)}
      -- the inner pipeline, also called directly by the multi-band
      path.
\end{itemize}

The inner pipeline computes:

\begin{align}
  \mathbf z_{\rm lab}(t) &= \mathsf{N}^{-1}\mathbf z_{\rm loops}(t),\\
  \mathbf z_\perp(t) &= \mathsf{P}_\perp \mathbf z_{\rm lab}(t),\\
  I(t) &= |\mathbf z_\perp|^2,\\
  A_p(t) &= \hat{\mathbf p}\cdot\mathbf z_\perp,\quad
  A_q(t) = \hat{\mathbf q}\cdot\mathbf z_\perp,
\end{align}
then forwards $(A_p, A_q)$ to \texttt{stokes.compute\_stokes()} which
returns the four Stokes parameters, the polarization fraction
$\sqrt{Q^2+U^2+V^2}/I$, the ellipticity $\frac12 \arcsin(V/I_p)$ in
degrees ($\pm 45^\circ$ at full circular polarization), and the
position angle $\frac12 \arctan2(U, Q)$.

A $2\times 2$ coherency on $(A_p, A_q)$ identifies the dominant-pol
eigenvector $u_{\rm dom}$, and
$z_{\rm dom}(t) = u_{\rm dom}^\top (A_p, A_q)$ provides an instantaneous
phase signal whose unwrapped, linearly-detrended residual gives both
the residual carrier frequency and a time-resolved frequency series via
\texttt{numpy.gradient}.  Everything is bundled into an
\texttt{AnalysisResult} dataclass.

\subsection{\texttt{triloop.multiband}}

Wraps band extraction and per-band analysis.  Public functions:

\begin{itemize}\itemsep=2pt
\item \texttt{analyze\_all\_bands(cap, az\_deg, el\_deg,
      loops\_channels, bw\_hz, decim\_rate\_hz, rf\_bands, loops\_config)}
      -- calls \texttt{extract\_bands} once and runs
      \texttt{analyze\_z\_loops} on each band.  Returns a dict mapping
      \texttt{rf\_hz} $\to$ \texttt{MultiBandResult} with fields
      \texttt{(rf\_hz, extraction, analysis, snrs\_db)}.
\item \texttt{make\_multiband\_figure(results, out\_path)} -- plot a
      five-row, four-column matplotlib figure (one row per band) of
      intensity, instantaneous frequency, ellipticity, and polarization
      fraction.
\end{itemize}

%-----------------------------------------------------------------------
\section{Direction finding}
%-----------------------------------------------------------------------

\subsection{\texttt{triloop.direction}}

Public functions, all operating on a stacked lab-frame complex array
$\mathsf{Z}_{\rm lab}\in\mathbb{C}^{3\times N}$:

\begin{itemize}\itemsep=2pt
\item \texttt{coherency\_matrix(Z\_lab)} returns
      $\mathsf{C} = \tfrac{1}{N}\mathsf{Z}_{\rm lab}\mathsf{Z}_{\rm lab}^H$,
      the $3\times 3$ Hermitian coherency.
\item \texttt{estimate\_direction\_from\_z\_lab(Z\_lab)} returns the
      smallest-eigenvalue eigenvector of $\mathsf{C}$ as
      $\hat{\mathbf k}_{\rm est}$, plus the polarization basis
      ($\hat{\mathbf p}$, $\hat{\mathbf q}$ from the two largest
      eigenvectors), the eigenvalues (descending), and a scalar
      \texttt{null\_strength} $\lambda_3/\lambda_1$.  This is the
      \emph{coarse} direction estimate; small \texttt{null\_strength}
      means the rank-1 trap (\S~\ref{sec:rank1ref}) is broken and
      the estimate is trustworthy.
\item \texttt{perp\_residual\_energy(Z\_lab, khat)} computes
      $\langle |\hat{\mathbf k}\cdot \mathbf z_{\rm lab}|^2 \rangle$,
      the quantity minimised by direction finding.
\item \texttt{null\_search(Z\_lab, az0, el0, half\_width\_deg, n\_points)}
      performs a 2-D grid sweep around $(az_0, el_0)$ and returns the
      minimiser plus the full residual map.
\item \texttt{parabolic\_refine(R, az\_grid, el\_grid)} fits a
      parabola to the $3\!\times\!3$ neighbourhood of the grid
      minimum to give a sub-grid direction.
\item \texttt{lock\_and\_analyze(...)} -- end-to-end entry point used
      by the \texttt{locate} CLI.
\end{itemize}

\subsection{Why find the null instead of the maximum?}
\label{sec:rank1ref}

The direction-formed power
$P(\hat{\mathbf k}) = \langle|\mathbf B_\perp(\hat{\mathbf k})|^2\rangle$
varies as $\frac12 + \frac12\cos^2\beta$ near the true direction
$\hat{\mathbf k}_{\rm true}$ (a shallow plateau).  The
\emph{perpendicular} residual
$Q(\hat{\mathbf k}) = \langle|\hat{\mathbf k}\cdot\mathbf B|^2\rangle$
varies as $\sin^2\beta\sim\beta^2$ near the same direction (a sharp
parabolic minimum).  Localising a parabolic minimum is far more
precise than a cosine-squared maximum, so we always lock on the null
rather than the peak.

For purely linearly polarized signals from a single direction,
$\mathsf{C}$ is rank 1 and the null is a 2-D plane rather than a
point.  In that regime \texttt{estimate\_direction\_from\_z\_lab}
returns \emph{a} direction (the smallest-eigenvalue one within
floating-point noise), but it is not physically meaningful; trustworthy
localisation requires Faraday rotation, a circularly polarized source,
or triangulation.

\subsection{\texttt{triloop.beamform}}

Older, simpler beamforming sweep used by some analysis notebooks.
Two functions:

\begin{itemize}\itemsep=2pt
\item \texttt{beamform\_grid(z\_loops, loops\_config, az\_grid\_deg,
      el\_grid\_deg)} -- compute time-averaged
      $P(\varphi, \theta) = \langle|\mathbf B_\perp(\hat{\mathbf k}_{\varphi,\theta})|^2\rangle$
      over a 2-D grid.
\item \texttt{best\_direction(P, az\_grid, el\_grid)} -- return the
      $\arg\max$ as $(\varphi, \theta)$.
\end{itemize}

\subsection{\texttt{triloop.stokes}}

A single function: \texttt{compute\_stokes(A\_p, A\_q)} returning a
dict with \texttt{I, Q, U, V, pol\_fraction, ellipticity\_deg,
position\_angle\_deg}.  Definitions:
$I = |A_p|^2 + |A_q|^2$,
$Q = |A_p|^2 - |A_q|^2$,
$U = 2\Re\{A_p^*A_q\}$,
$V = 2\Im\{A_p^*A_q\}$.

%-----------------------------------------------------------------------
\section{Reporting and visualization}
%-----------------------------------------------------------------------

\subsection{\texttt{triloop.view}}

Static QC PNG generator used by \texttt{triloop view}.
\texttt{make\_view\_figure(cap, rf\_bands, bw\_hz, out\_path)}
produces a single matplotlib figure with three rows: a full-Nyquist
PSD with each RF band annotated at its alias position, per-band PSD
zooms (with the displayed Δf-axis flipped for inverted bands so the
operator reads frequencies the right way), a decimated time-domain
trace of every channel, and the auto-range probe table.

\subsection{\texttt{triloop.browse}}

Interactive single-file HTML browser.  Top-level entry points:

\begin{itemize}\itemsep=2pt
\item \texttt{build\_browser\_html(cap, az\_deg, el\_deg,
      loops\_channels, bw\_hz, decim\_rate\_hz, n\_anim\_frames,
      title)} -- returns a self-contained HTML string.
\item \texttt{write\_browser(cap, out\_path, **kw)} -- save it to disk.
\end{itemize}

The page uses Plotly (loaded once via CDN; if you need a fully offline
build, change the \texttt{include\_plotlyjs="cdn"} argument to
\texttt{True} and the page will inline the entire library at the cost
of $\sim 3$\,MB extra).  Layout: tab bar at the top with
\textit{Overview}, \textit{Time series}, and one tab per RF band.
The per-band tab has four panels: PSD zoom, intensity time history
(linear + togglable dB), intensity CDF with P10/P1 fade-depth markers,
and an animated polarization ellipse with Play/Pause and a frame
slider.

The animation is implemented as Plotly frames; each frame draws a
sliding window of $(\Re A_p, \Re A_q)$ samples centred on a different
slow-time epoch.  The window length is chosen automatically to be a
few cycles of the residual carrier-phase variation at the decimated
rate.

%-----------------------------------------------------------------------
\section{Module roll call}
%-----------------------------------------------------------------------

This section lists every Python file in both packages and what each
one provides.  Treat it as a table of contents for the source.

\subsection{\textsf{picoacq}}

\begin{tabular}{p{0.21\linewidth}p{0.74\linewidth}}
\toprule
\textbf{Module} & \textbf{Purpose} \\\midrule
\texttt{\_\_init\_\_.py} & package marker; no public API of its own. \\
\texttt{cli.py}   & Click-based CLI; subcommands \texttt{capture} and
                    \texttt{monitor}.  Reads
                    \texttt{--rate / --duration / --channels /
                    --range / --auto-range / --headroom / --probe /
                    --rf-bands / --simulate / --sim-*} options and
                    forwards them to \texttt{recorder.capture()}. \\
\texttt{recorder.py} & Top-level \texttt{capture()} orchestrator with
                    \texttt{\_check\_buffer()} guard,
                    \texttt{auto\_range()} probe, range chooser
                    \texttt{\_pick\_range()}, and the
                    \texttt{\_try\_real\_capture()} hardware path
                    (with simulator fallback). \\
\texttt{ps5444a.py} & PicoSDK 5000a wrapper: \texttt{open\_5444()},
                    \texttt{configure\_channel()},
                    \texttt{get\_timebase()}, \texttt{acquire\_block()}.
                    Plus \texttt{\_ensure\_pico\_dyld\_path()} for
                    macOS framework search. \\
\texttt{simulator.py} & \texttt{simulate\_capture()} -- 3-loop
                    no-hardware fallback compatible with the
                    \texttt{recorder} return shape. \\
\texttt{alias.py}    & RF~$\leftrightarrow$~baseband mapping under
                    bandpass undersampling.  Public:
                    \texttt{Alias} dataclass,
                    \texttt{alias\_of()},
                    \texttt{aliases\_for\_bands()},
                    \texttt{rf\_for\_baseband()}.  See
                    Section~\ref{sec:alias}. \\
\bottomrule
\end{tabular}

\subsection{\textsf{triloop}}

\begin{tabular}{p{0.21\linewidth}p{0.74\linewidth}}
\toprule
\textbf{Module} & \textbf{Purpose} \\\midrule
\texttt{\_\_init\_\_.py} & re-exports the public API (see
                    \textit{User Manual} for the full list). \\
\texttt{cli.py}   & Click-based CLI; subcommands
                    \texttt{view, analyze-multi, browse, analyze,
                    locate, batch, summary, simulate, notebook}. \\
\texttt{io\_hdf5.py}  & \texttt{read\_capture()} and
                    \texttt{write\_capture()} for the HDF5 v0.1 layout
                    described in \S\ref{sec:fileformat-ref}. \\
\texttt{config.py}    & \texttt{default\_loops\_config()} (cube-vertex
                    geometry) and \texttt{validate\_loops\_config()}
                    schema check. \\
\texttt{geometry.py}  & \texttt{az\_el\_to\_khat()},
                    \texttt{build\_N\_matrix()},
                    \texttt{perp\_projector()},
                    \texttt{perp\_orthonormal\_basis()},
                    \texttt{LOOP\_NORMALS\_DEFAULT}. \\
\texttt{extract.py}   & Single-band complex-baseband extractor:
                    \texttt{extract\_complex\_baseband()},
                    \texttt{extract\_three\_loops()}. \\
\texttt{bands.py}     & Multi-band extractor:
                    \texttt{extract\_bands()},
                    \texttt{BandExtraction} dataclass. \\
\texttt{stokes.py}    & \texttt{compute\_stokes(A\_p, A\_q)} ->
                    Stokes parameters and derived quantities. \\
\texttt{analyze.py}   & Top-level \texttt{analyze()} and inner
                    \texttt{analyze\_z\_loops()}; \texttt{AnalysisResult}
                    dataclass. \\
\texttt{multiband.py} & \texttt{analyze\_all\_bands()} (extract +
                    analyze for every band) and
                    \texttt{make\_multiband\_figure()}. \\
\texttt{direction.py} & Coherency-based direction finding:
                    \texttt{coherency\_matrix()},
                    \texttt{estimate\_direction\_from\_z\_lab()},
                    \texttt{perp\_residual\_energy()},
                    \texttt{null\_search()},
                    \texttt{parabolic\_refine()},
                    \texttt{lock\_and\_analyze()}. \\
\texttt{beamform.py}  & Older/simpler grid-sweep:
                    \texttt{beamform\_grid()},
                    \texttt{best\_direction()}. \\
\texttt{view.py}      & Static QC PNG: \texttt{make\_view\_figure()}. \\
\texttt{browse.py}    & Interactive Plotly HTML browser:
                    \texttt{build\_browser\_html()},
                    \texttt{write\_browser()}, plus the helpers
                    \texttt{\_channel\_psd()},
                    \texttt{\_band\_zoom\_spectrum()},
                    \texttt{\_smooth\_decimate()},
                    \texttt{\_cdf()},
                    \texttt{\_build\_pol\_animation()}. \\
\bottomrule
\end{tabular}

%-----------------------------------------------------------------------
\section{Tracing a capture through the code}
%-----------------------------------------------------------------------

To consolidate the picture above, here are the call paths for the
three new commands:

\paragraph{\texttt{triloop view cap.h5}.}
\begin{enumerate}\itemsep=2pt
\item \texttt{cli.view\_cmd} parses options.
\item \texttt{io\_hdf5.read\_capture(cap.h5)} -> dict.
\item \texttt{view.make\_view\_figure(cap, rf\_bands, bw\_hz, out\_path)}.
\item Inside \texttt{make\_view\_figure}: per-channel PSD via
      \texttt{\_channel\_psd}; per-band zoom via
      \texttt{\_band\_zoom\_spectrum} (calls
      \texttt{picoacq.alias.alias\_of} for each band centre);
      time-domain decimated; probe table from
      \texttt{capture\_settings.auto\_range\_probe}.
\end{enumerate}

\paragraph{\texttt{triloop analyze-multi cap.h5 --az 9 --el 35}.}
\begin{enumerate}\itemsep=2pt
\item \texttt{cli.analyze\_multi\_cmd} parses options.
\item \texttt{io\_hdf5.read\_capture(cap.h5)} -> dict.
\item \texttt{multiband.analyze\_all\_bands(cap, az, el, ...)}:
\begin{enumerate}\itemsep=1pt
\item \texttt{bands.extract\_bands(cap, ...)} which for each band uses
      \texttt{picoacq.alias.alias\_of} to find $f_{\rm bb}$, then mixes
      and decimates each channel via \texttt{\_extract\_one\_channel}.
\item For each band: \texttt{\_per\_loop\_snr\_from\_baseband}, then
      \texttt{analyze.analyze\_z\_loops(t, Z\_loops, f\_rf, ...)}.
\end{enumerate}
\item \texttt{multiband.make\_multiband\_figure(results, out\_png)}.
\item Per-band JSON summary written to \texttt{out\_json}.
\end{enumerate}

\paragraph{\texttt{triloop browse cap.h5}.}
\begin{enumerate}\itemsep=2pt
\item \texttt{cli.browse\_cmd} parses options.
\item \texttt{io\_hdf5.read\_capture(cap.h5)} -> dict.
\item \texttt{browse.write\_browser(cap, out\_html, ...)} which calls
      \texttt{build\_browser\_html(cap, ...)}:
\begin{enumerate}\itemsep=1pt
\item Calls \texttt{multiband.analyze\_all\_bands} (same path as
      above) to populate per-band intensity, ellipticity, and
      $A_p / A_q$ time series.
\item Builds per-tab Plotly figures: overview PSD, time series, and
      one 4-panel tab per band.
\item For each band assembles a polarization-ellipse animation via
      \texttt{\_build\_pol\_animation()} (Plotly frames + slider).
\item Renders the page using \texttt{plotly.io.to\_html()} for each
      figure into a tabbed HTML scaffold.
\end{enumerate}
\end{enumerate}

%-----------------------------------------------------------------------
\section{Testing and verification}
%-----------------------------------------------------------------------

The test suite under \texttt{triloop/tests/} exercises the analysis
pipeline against the simulator at known polarizations and directions.
Useful commands during development:

\begin{lstlisting}[style=shell]
# Create a synthetic 4-channel capture with all 5 WWV bands aliased
PYTHONPATH=triloop:picoacq python3 -m picoacq.cli capture \
    --simulate --duration 0.5 -o sim.h5

# Run the multi-band analysis and confirm pol_frac and ellipticity
PYTHONPATH=triloop:picoacq python3 -m triloop.cli analyze-multi sim.h5 \
    --az 9 --el 35

# Render the interactive browser
PYTHONPATH=triloop:picoacq python3 -m triloop.cli browse sim.h5
\end{lstlisting}

For the end-to-end ``does the alias arithmetic come out right'' check
on a synthetic 4-channel 15.625\,MS/s capture, see
\texttt{triloop/triloop/bands.py} (specifically \texttt{extract\_bands}
will raise \texttt{ValueError} if the requested per-band bandwidth
overlaps the Nyquist edge or another band's alias, which is the
nearest thing to an automatic correctness check).

\end{document}
